\section{Introduction}

2024 marked the first calendar year that the global mean temperature exceeded 1.5 degrees above pre-industrial levels. This means that the lower bound maximum temperature set out in the Paris Climate Agreement has been reached \cite{worldmeteorologicalorganizationWMOConfirms20242025}. % TODO: Add Paris Agreement citation to references.bib (entry #2 from lut.txt: United Nations Framework Convention on Climate Change. (2016). The Paris Agreement. UNFCCC. https://unfccc.int/sites/default/files/resource/parisagreement_publication.pdf) The impact of this warming will be particularly pronounced in urban areas, where more than half of the world's population currently resides \cite{SecretaryGeneralReport2023}, and where temperatures are typically 2-5°C higher when compared to lower-density areas, due to the Urban Heat Island effect \cite{okeUrbanClimates2017a,mohajeraniUrbanHeatIsland2017a}. This temperature differential results primarily from the altered energy balance caused by differences in surface materials between urbanized and rural areas \cite{okeEnergeticBasisUrban1982a,santamourisAnalyzingHeatIsland2015a}. UHI has been documented to negatively affect human health through increased heat-related mortality \cite{heavisideUrbanHeatIsland2017a,watts2018ReportLancet2018,faurieAssociationHighTemperature2022}, with studies showing 2-4\% productivity decreases per degree Celsius of air temperature above comfort thresholds \cite{kjellstromHeatHumanPerformance2016,speakInfluenceTreeTraits2020,rahmanTreeCoolingEffects2020}. Therefore, strategies that counter UHI effects can significantly contribute to improving urban livability. Urban trees represent one such effective strategy, with research demonstrating their dual capacity to mitigate UHI by altering surface energy balance: reducing ambient temperatures by 0.5-2.0°C in their immediate vicinity \cite{rahmanTraitsTreesCooling2020a,liCoolingEfficacyTrees2024}, as well as helping urban areas adapt by providing direct shade for inhabitants, which can reduce mean radiant temperatures by up to 20°C \cite{duFieldAssessmentNeighboring2020,middelImpactShadeOutdoor2016,liCoolingEfficacyTrees2024}.

Traditional methods for urban forest management often rely on manual surveys conducted by trained arborists. These labor-intensive processes typically require 20-30 minutes per tree \cite{nielsenReviewUrbanTree2014}, resulting in infrequent updates and limited spatial coverage at the city scale. This approach yields low spatial and temporal granularity of tree data. The City of Amsterdam has made progress by maintaining a comprehensive tree database with information on over 260,000 public trees. However, the city estimates the total number of trees in Amsterdam ranges from 800,000 to 1,000,000 when private trees are included \cite{cityofamsterdamAmsterdamGreenInfrastructure2020}, indicating that approximately 70\% of the urban forest remains undocumented. In the existing tree census, key attributes such as tree height and age are often provided in broad ranges rather than precise measurements, further limiting their utility for climate modeling. As urban areas expand and urban forests grow, there is an increasing need for more comprehensive and accurate information to assess their effect on urban climate and to optimize maintenance resources. Recent studies \cite{nowakUSUrbanForest2018a,thomasTreesOutsideForests2021,przewoznaGeoQuestionnaireEnvironmentalPlanning2021,konijnendijkDefiningUrbanForestry2006} demonstrate that incomplete tree inventories can lead to underestimations of ecosystem services by 15-35\%, highlighting the practical importance of improved assessment methods.
 
To address these limitations, we incorporate individual-tree Leaf Area Index (LAI) assessment in our methodology. While LAI is widely used in ecological studies, its application within urban forestry, especially for individual trees, remains underexplored. LAI is a metric quantifying canopy density, defined as the total one-sided leaf area per unit of ground surface area, expressed as a dimensionless ratio (m²/m²) \cite{chenDefiningLeafArea1992}. This metric provides critical insights into an urban tree's capacity to intercept sunlight. Higher LAI values indicate greater leaf coverage, which increases shading capacity. Research by \cite{shashua-barMicroclimateModellingStreet2010} demonstrates that trees with LAI values above 3 can reduce surface temperatures by 8-12°C compared to 3-5°C for trees with lower LAI values. Furthermore, higher LAI correlates with increased photosynthetic capacity and biomass production \cite{reichKeyCanopyTraits2012}. The ecological functions of urban trees quantified by this measure extend to carbon sequestration, with forests having higher LAI demonstrating 25-40\% greater potential for carbon storage \cite{hardimanCouplingFineScaleRoot2017}. In urban contexts, LAI also helps quantify benefits of trees in terms of air quality improvement, with each unit increase in LAI associated with 1-4\% reduction in particulate matter concentration \cite{nowakTreeForestEffects2014}, stormwater management benefits of 2-7\% increased retention per LAI unit \cite{berlandRoleTreesUrban2017a}, and improved habitat quality for various urban wildlife species \cite{threlfallApproachesUrbanVegetation2016}. By integrating LAI estimation, our approach provides more comprehensive insights into shade mapping and the ecosystem benefits of trees in cities.

This paper proposes an automated approach for urban forestry management that addresses the limitations of traditional inventory methods. Our methodology employs a voxelization technique (a three-dimensional discretization of space into volumetric pixels) for LAI estimation and leverages deep learning segmentation models to process and analyze three complementary aerial data sources: aerial LiDAR scans with a density of 25 points/m², high-resolution (25cm) RGB orthophotos, and near-infrared imagery. We integrate these data sources to detect climate-relevant morphological features at an individual tree level with significantly higher precision than conventional methods. The approach focuses on extracting key parameters such as tree height, canopy diameter, and LAI, which previous research has identified as the most crucial variables for estimating shading potential and overall ecosystem services \cite{rahmanTraitsTreesCooling2020a,threlfallApproachesUrbanVegetation2016}. The final component of our methodology is a machine learning regression model that predicts LAI from morphological data and taxonomic classification, enabling the scaling of our approach to regions where high-resolution LiDAR data may not be available.

The significance of this research is its potential to enhance urban forest management by making tree inventories more accessible, accurate, and actionable for urban planners. Our approach increases accessibility through: (1) automating labor-intensive field measurements that traditionally require specialized arborist expertise; (2) substantially reducing costs associated with manual inventories through remote sensing techniques; and (3) creating standardized methodologies that can be implemented by municipalities with varying technical capacities using widely available remote sensing data. Automating the inventory process and including private property trees provides a comprehensive view of the urban forest and its climate mitigation potential. Integrating LAI mapping offers a quantitative basis for shade assessment, aiding strategic tree planting in heat-vulnerable areas. As cities face thermal challenges due to climate change, tools enabling evidence-based urban forestry decisions are crucial. Our approach bridges remote sensing technology with practical urban planning, contributing to more climate-resilient urban environments. The following sections detail our methodology, present results from Amsterdam, and discuss implications for future urban forestry practices.
